\section{Conclusion}
\label{sec:conclusion}

Kernel Search for the large-scale MSCFLP is effective because it uses the LP
relaxation to identify a small set of promising facilities and arcs before
solving restricted MILPs.
The same dependence on the LP, however, exposes a bottleneck on the largest
instances: obtaining the LP information can itself become expensive.
This paper addresses that bottleneck directly by solving the KS LP relaxation
through column generation.
The intention is not to alter the heuristic guidance, but to obtain the
LP-derived bound, facility partition, and reduced-cost signal at lower setup
cost.

The second contribution concerns the restricted-MILP phase.
The original KS framework already contains several speed controls, such as
limiting bucket exploration and removing unpromising facilities from the kernel.
Those controls reduce MILP effort but naturally introduce a speed-quality
trade-off.
FKS instead organizes the restricted search as a flat-$k$ funnel.
The flat-$k$ rule gives predictable base subproblem sizes, while preserving all
incumbent arcs and facilities between consecutive stages guarantees that the
previous solution is feasible in the next restricted MILP.
This gives a simple structural basis for warm-starting later stages after the
edge set has been narrowed.

The broader message is diagnostic.
Matheuristic improvements benefit from examining the existing method's design
choices carefully before proposing new components.
For G\&S~2012 KS on the MSCFLP, a close reading surfaces three design aspects
that can each be refined --- avoidable LP cost, biased arc selection, and
cold-start MILP transitions --- each with a targeted, independently verifiable
modification.
The computational study is organized around this structure: each comparison
changes exactly one component, making it possible to attribute improvements
to a specific design decision rather than to the aggregate effect of the method.

Several extensions remain natural.
First, completing the remaining Test Bed~A and Test Bed~C families will provide
a broader robustness check across facility-client scales and capacity regimes.
Second, per-client adaptive budgets could allocate more arcs to clients with
diffuse LP support and fewer arcs to clients with concentrated support while
retaining the incumbent-preservation guarantee.
Third, the feasibility-preservation principle may be useful in other Kernel
Search settings whenever restricted subproblems are built by selecting subsets
of variables used by an incumbent solution.
